\SourcePage{677}
\SGAchapter{XIV}{相交理论中的未解决问题}{Problèmes ouverts en théorie des intersections}
\begin{center}A. Grothendieck\end{center}
本讲回顾与本讨论班所遇问题有关、值得研究的若干问题。我们不讨论 Serre 意义下的\Term{local-scheme}{局部}\Term{intersection-theory}{相交理论} [14] 所涉及的问题，因为本讨论班没有用到这一理论；关于这个方向上仍未解决的两个基本猜想，特别请参见该书第 V 章第 4 节。我们也不谈那些对算术应用十分重要、涉及\Term{cycle-class-group}{圈类群}结构的问题，尤其是\Term{smooth}{光滑}\Term{projective}{射影}\Term{algebra}{代数}\Term{scheme}{概形}上的这类问题；对此请参阅 S. Kleiman 在《\Term{scheme}{概形}\Term{cohomology}{上同调}十讲》（\textit{Dix exposés sur la cohomologie des schémas}，Masson--North Holland Publishing Cie）中的《\Term{algebraic-cycle}{代数圈}与 Weil 猜想》（\textit{Algebraic cycles and the Weil conjectures}）。本讲第 4 节还给出本讨论班的一些补充，尤其是有关 \Term{chow-ring}{Chow 环}的补充；这些内容在前面的各讲中没有安排进去。

\SGAsection{1}{导出范畴 $D(X)$ 中的 $\Lambda$ 运算}
\footnote{参见原讲第 23 页的补注（本译文 PDF 原页 699）。}
设 $X$ 是\Term{locally-ringed-topos}{局部赋环拓扑斯}，$K'(X)$ 是 $X$ 上\Term{perfect-complex}{完美复形}的类所成的环（第 IV 讲）。我们希望在 $K'(X)$ 上定义一个关于 $X$ 具有\Term{functorial}{函子性}的 $\lambda$-\Term{lambda-ring}{环}结构（第 V 讲），使自然\Term{homomorphism}{同态}
\[K'_{\text{naïf}}(X)\longrightarrow K'(X)\]
成为一个\SGAsourcegap{原页此处“un”与“-homomorphisme”之间缺少修饰\Term{signature}{符号}}\Term{homomorphism}{同态}；这里左边是 $X$ 上\Term{locally-free}{局部自由}\Term{module}{模}的类所成的环，配备第 V 讲定义的 $\lambda$ 结构。为此，应定义“\Term{exterior-power}{外幂}”\Term{functor}{函子}（\Term{generic}{一般}并非\Term{additive}{加性}）
\[\bigwedge^i:D^-(X)\longrightarrow D^-(X),\]
使其将\Term{pseudo-coherent}{伪凝聚}\Term{complex}{复形}（分别为\Term{perfect-complex}{完美复形}）变为\Term{pseudo-coherent}{伪凝聚}\Term{complex}{复形}（分别为\Term{perfect-complex}{完美复形}），并且对于\Term{exact-triangle}{正合三角}满足必要的\Term{form}{形式}性质，从而能用下式在 $K'(X)$ 上定义 $\lambda$ 结构：
\[\lambda^i\bigl(\operatorname{cl}(L^\bullet)\bigr)
=\operatorname{cl}\bigl(\bigwedge^i L^\bullet\bigr).\tag{1.1}\]
\SourcePage{678}
当 $X$ 的\Term{characteristic}{特征}为零，即其\Term{structure-sheaf}{结构层}是 $\mathbf Q$-\Term{algebra}{代数}\Term{sheaf}{层}时，只需取
\[\bigwedge^i L^\bullet
=\left(\mathop{\bigotimes}\nolimits^{i,\mathbf L}L^\bullet\right)^{\mathrm{alt}},\tag{1.2}\]
问题就得到解决。这里把 $L^\bullet$ 的第 $i$ 个\Term{derived-category}{导出范畴}意义下的\Term{tensor-power}{张量幂}看成 $D^-(\mathscr O_i)$ 中的一个\Term{object}{对象}；$\mathscr O_i$ 是\Term{symmetric-group}{对称群} $\mathfrak S_i$ 在 $\mathscr O_X$ 上的\Term{group-algebra}{群代数}，上标 $\mathrm{alt}$ 表示取“$\mathfrak S_i$ \Term{action}{作用}下的交替部分”（\Term{characteristic}{特征}为零的假设使这一运算有意义）。在 [8]（下文记为 [RRR]）中，为首次证明 \Term{riemann-roch-theorem}{Riemann--Roch 定理}，已经用过这类公式。

不作\Term{characteristic}{特征}假设时，若 $L^\bullet$ 在\Term{degree}{次数} $>0$ 处为零，则根据 Dold--Puppe [6]，$L^\bullet$ \Term{canonical}{典范}地\Term{homotopic}{同伦}于一个\Term{semi-simplicial-module}{半单纯模} $L^{\bullet\prime}$ 所对应的\Term{complex}{复形}，而 $L^{\bullet\prime}$ 本身也由 $L^\bullet$ \Term{canonical}{典范}地给出。因此，对于 $\bigwedge^i L^\bullet$（例如设 $L^\bullet$ \Term{flat}{平坦}），一个颇自然的候选是\Term{semi-simplicial-module}{半单纯模} $\bigwedge^i L^{\bullet\prime}$ 对应的\Term{complex}{复形}，其\Term{degree}{次数} $n$ 的\Term{set-component}{分量}为 $\bigwedge^i L'_n$。但要把定义推广到 $D^-(X)$ 中不再假定在\Term{degree}{次数} $>0$ 处为零的\Term{complex}{复形}，就必须研究上述运算与\Term{translation}{平移}运算 $L^\bullet\mapsto L^\bullet[1]$（按 [6] 的术语称为\Term{suspension}{悬挂}）之间的关系。遗憾的是，似乎并不存在\Term{isomorphism}{同构}
\[\bigwedge^i\bigl(L^\bullet[1]\bigr)\simeq
\bigl(\bigwedge^i L^\bullet\bigr)[i],\]
因而不能借此直接完成所希望的推广。

\SGAsection{2}{不作射影性假设的 Riemann--Roch 公式}
\textbf{2.1。}一旦各个 $K'(X)$ 上有了自然的 $\lambda$ 结构，就可以按第 VIII 讲所考虑的\Term{form}{形式}，对任意\Term{proper}{固有}\Term{complete-intersection-morphism}{完全交态射} $f:X\to Y$ 表述 \Term{riemann-roch-theorem}{Riemann--Roch 定理}；甚至更\Term{generic}{一般}地，可考虑\Term{locally-ringed-topos}{局部赋环拓扑斯} $Y$ 上的相对\Term{scheme}{概形} $X$，而 $X$ 仍满足相应的相对条件。$f$ 的\Term{todd-class}{Todd 类}可由第 0 讲 4.4 构造的 $f$ 的\Term{relative-cotangent-complex}{相对余切复形} $T_f$ 在 $K'(X)$ 中的类来定义；由于对 $f$ 所作的\Term{local-scheme}{局部}假设，这个\Term{complex}{复形}在此是\Term{perfect-complex}{完美的}。显然，在这一\Term{generic}{一般}情形下证明 \Term{riemann-roch-formula}{Riemann--Roch 公式}，需要引入本质上全新的思想；这些思想也可能提供在复解析情形
\SourcePage{679}
（甚至\Term{rigid-analytic}{刚性解析}情形！）下证明类似公式的关键。回顾一下：在复解析情形，目前已知的只有\Term{nonsingular}{非奇异}\Term{compact}{紧}簇 $X$ 上\Term{locally-free}{局部自由}\Term{sheaf}{层} $F$ 的情况，且底\Term{space}{空间} $Y$ 为一点（这归功于 Atiyah--Singer 指标公式）；若只假定 $F$ \Term{coherent}{凝聚}，现有技术似乎仍无能为力。

\textbf{2.2。}作为 \Term{riemann-roch-formula}{Riemann--Roch 公式}的一个基本组成部分，也再提醒读者注意第 III 讲提出的\Term{finiteness-problem}{有限性问题}：若 $f:X\to Y$ 是\Term{proper}{固有}\Term{morphism}{态射}，而 $L^\bullet$ 是 $X$ 上相对于 $Y$ \Term{pseudo-coherent}{伪凝聚}的\Term{complex}{复形}，则 $\mathrm Rf_*(L^\bullet)$ 是否\Term{pseudo-coherent}{伪凝聚}？这个猜想的\Term{generic}{一般}证明很可能会运用一些\Term{algebra}{代数}思想，并由此在比原文宽泛得多的条件下证明 Grauert 的\Term{finiteness}{有限性}“\Term{theorem}{定理}”[7]：从一开始就允许底\Term{space}{空间}（或\Term{topos}{拓扑斯}）$Y$ 是带有\Term{topology}{拓扑} $\mathbf C$-\Term{algebraic-structure}{代数结构}的\Term{ringed-space}{赋环空间}，范围比\Term{analytic-space}{解析空间}更广，包括 $Y$ 为配备\Term{continuous}{连续}复值函数\Term{sheaf}{层}的\Term{locally-compact}{局部紧}\Term{space}{空间}，或配备 $C^m$ 复值函数\Term{sheaf}{层}的 $C^m$ \Term{manifold}{流形}，或\Term{real-analytic-space}{实解析空间}等情形。

\textbf{2.3。}最后指出，在这些问题中，应当把 $f$ 的\Term{properness}{固有性}假设替换为这样的假设：所考虑的 $X$ 上\Term{complex}{复形}的\Term{cohomology-sheaf}{上同调层} $H^i(L^\bullet)$ 的\Term{support}{支撑}在 $Y$ 上\Term{proper}{固有}。\footnote{译注：本句原文印作“aient leurs groupes propres sur $Y$”；这里的“\Term{support}{支撑}”据本句的数学语义译出，原文的 groupes 一词保留记录于报告。}

\SGAsection{3}{浸入的“无分母”Riemann--Roch 公式}
\textbf{3.1。}设 $f:Y\to X$ 是\Term{regular-closed-immersion}{正则闭浸入}（第 VII 讲），$N$ 为其\Term{conormal-sheaf}{余法层}。为简单起见，假定 $X$ 有一个\Term{ample}{丰沛}\Term{invertible-sheaf}{可逆层}。对 $y\in K'(Y)$ 的 \Term{riemann-roch-formula}{Riemann--Roch 公式}
\[\operatorname{ch}_X(f_*(y))
=f_*\bigl(\operatorname{ch}_Y(y)\operatorname{Todd}(-\check N)\bigr),\]
结合\Term{generic}{一般}公式（第 VII 讲）
\[f_*(y')f_*(y'')=f_*\bigl(y'y''c_d(\check N)\bigr),
\qquad y',y''\in\operatorname{Gr}(Y)_{\mathbf Q}\]
（这里假定 $f$ 的\Term{codimension}{余维数}为 $d$），可以把 $f_*(y)$ 在 $\operatorname{Gr}^\bullet(X)_{\mathbf Q}$ 中的\Term{chern-class}{Chern 类}写成
\SourcePage{680}
\[c_i(f_*(y))
=f_*\left(P_i\bigl((c_\alpha(y))_\alpha,(c_\beta(\check N))_\beta\bigr)\right),\tag{3.1}\]
其中 $P_i$ 是关于 $c_\alpha(y)$、$c_\beta(\check N)$ 的\Term{universal-polynomial}{通用多项式}，其系数先验地是有理数。事实上，可以发现 [RRR]，这些\Term{polynomial}{多项式}具有整数系数。这就提出了“无分母”\Term{riemann-roch-formula}{Riemann--Roch 公式}的\Term{effectivity}{有效性}问题。实际上，对于 $c_i(f_*(y))$，\Term{generic}{一般}不能期望在 $\operatorname{Gr}^\bullet(X)$ 本身中有这种公式，因为\Term{homomorphism}{同态} $f_*:K'(Y)\to K'(X)$ 与 $\lambda$-\Term{filtration}{滤过}（相差 $d$ 的位移）本身并不相容（4.6），只有模去\Term{torsion}{挠元}才相容。因此，不能定义一个\Term{homomorphism}{同态}
\[f_*:\operatorname{Gr}^\bullet(Y)\longrightarrow\operatorname{Gr}^\bullet(X),\]
使它给出自然\Term{homomorphism}{同态} $f_*:\operatorname{Gr}^\bullet(Y)_{\mathbf Q}\to\operatorname{Gr}^\bullet(X)_{\mathbf Q}$，从而把公式 (3.1) 理解为 $\operatorname{Gr}^\bullet(X)$ 中、而不只是在 $\operatorname{Gr}^\bullet(X)_{\mathbf Q}$ 中的公式。但若假定 $X$ \Term{regular}{正则}，因而 $K'(X)\simeq K_\bullet(X)$（第 IV 讲 2.5），则可在 $K'(X)$ 上赋予按\Term{support}{支撑}的\Term{codimension}{余维数}定义的\Term{topological-filtration}{拓扑滤过}，它继承自 $K_\bullet(X)$ 的相应\Term{filtration}{滤过}。可以合理地预期，这个\Term{filtration}{滤过}与\Term{ring-structure}{环结构}相容；至少当 $X$ 在域 $k$ 上\Term{smooth}{光滑}时，利用 \Term{chow-ring}{Chow 环}理论 [4] 可以证明这一点 [RRR]。此时，其\Term{associated-graded}{伴随分次对象}
\[\operatorname{Gr}_{\mathrm{top}}^\bullet(X)\]
就是\Term{graded-ring}{分次环}；再假定 $Y$ 也\Term{regular}{正则}，就具备了定义一个\Term{degree}{次数}为 $d$ 的\Term{graded-group}{分次群}自然\Term{homomorphism}{同态}所需的一切：
\[f_*=f_*^{\mathrm{top}}:
\operatorname{Gr}_{\mathrm{top}}^\bullet(Y)\longrightarrow
\operatorname{Gr}_{\mathrm{top}}^\bullet(X).\]
由于 $K'(X)$ 的 $\lambda$-\Term{filtration}{滤过}比刚才按\Term{codimension}{余维数}考虑的\Term{topological-filtration}{拓扑滤过}更细（第 X 讲 1.3.3），便得到\Term{graded-group}{分次群}（甚至\Term{graded-ring}{分次环}）的\Term{canonical}{典范}\Term{homomorphism}{同态}
\[\operatorname{Gr}^\bullet(X)\longrightarrow\operatorname{Gr}_{\mathrm{top}}^\bullet(X),\]
它在模去\Term{torsion}{挠元}后还是\Term{isomorphism}{同构}（第 VIII 讲\SGAsourcegap{原页“Exp. VIII”后未印节号}）。这就允许我们像 [RRR] 那样，考虑 $X$ 上取值于 $\operatorname{Gr}_{\mathrm{top}}^\bullet(X)$ 的\Term{chern-class}{Chern 类}理论，对 $Y$ 也一样，从而把公式 (3.1) 理解为 $\operatorname{Gr}_{\mathrm{top}}^\bullet(X)$ 中的公式。在这种\Term{form}{形式}下，[RRR] 已在 $X$ 的\Term{characteristic}{特征}为零时证明此公式，至少当 $X$ 在某个域上\Term{smooth}{光滑}时如此（这个假设仅用于保证 $K'(X)$ 的\Term{topological-filtration}{拓扑滤过}与\Term{ring-structure}{环结构}相容）。\Term{characteristic}{特征}为零的假设主要用于
\SourcePage{681}
取得 (1.1)、(1.2) 一类的公式；一旦第 1 节的问题得到肯定解决，本质上同一个证明应当不需\Term{characteristic}{特征}假设就能适用。

\textbf{3.2。}当 $X$ 在域上\Term{smooth}{光滑}时，还有 \Term{chow-ring}{Chow 环}
\[A(X)\]
以及 $X$ 上取值于 $A(X)$ 的\Term{chern-class}{Chern 类}理论 [9]，它比刚引入的取值于 $\operatorname{Gr}_{\mathrm{top}}^\bullet(X)$ 的理论更细：后者的\Term{chern-class}{Chern 类}是 Chow \Term{form}{形式}的\Term{chern-class}{Chern 类}在\Term{canonical}{典范}\Term{surjective}{满}\Term{homomorphism}{同态}（参见 4.1）
\[\phi:A(X)\longrightarrow\operatorname{Gr}_{\mathrm{top}}^\bullet(X)\]
下的像。这个\Term{homomorphism}{同态}模去\Term{torsion}{挠元}后是\Term{isomorphism}{同构}，但\Term{generic}{一般}并非\Term{isomorphism}{同构}（4.7）。在此前提下，若 $X,Y$ 均\Term{smooth}{光滑}，包含 $f:Y\to X$ 还定义一个\Term{degree}{次数}为 $d$ 的\Term{graded-group}{分次群}\Term{homomorphism}{同态}
\[f_*:A(Y)\longrightarrow A(X),\]
从而可把 (3.1) 理解为 $A(X)$ 中的公式。\footnote{这个公式以及 3.3 指出的变体，此后已由 P. Jouanolou 证明（1969 年 11 月）。参见\SGAsourcegap{原脚注末尾“Cf.”之后未给出引文}。}
当 $i<d$ 时，此公式\Term{degeneration}{退化}为 $0=0$，因此只有 $i\geq d$ 时才有意义。即使在 $i=d=2$ 时，这个公式目前也尚未证明（参见 4.3）；把\Term{algebra}{代数}\Term{cycle}{圈}的\Term{rational-equivalence}{有理等价}换成\Term{algebraic-equivalence}{代数等价}所得到的较弱公式，也尚未证明。与 $\operatorname{Gr}_{\mathrm{top}}^\bullet(X)$ 中的类似公式不同，似乎没有很有说服力的理由，足以使人相信无分母公式在 $A(X)$ 中必然成立。鉴于我们对 $i\geq2$ 时各群 $A^i(X)$ 的性质几乎一无所知，一个较好的检验办法似乎是：检验由它推\Term{derive}{导出}的 $A^1(Z)=\operatorname{Pic}(Z)$ 中的公式。例如，用 $g_*:A(X)\to A(Z)$ \Term{action}{作用}于 (3.1) 两边，其中 $g:X\to Z$ 是到一个\Term{dimension}{维数}为 $\dim X-i+1$ 的\Term{smooth}{光滑}\Term{quasi-projective}{拟射影}\Term{scheme}{概形} $Z$ 的\Term{proper}{固有}\Term{morphism}{态射}。

\textbf{3.3。}最后，还可利用\Term{chern-class}{Chern 类}的\Term{cohomology}{上同调}理论（SGA 5，第 VII 讲）赋予 (3.1) 意义，只要有 \Term{gysin-homomorphism}{Gysin 同态}
\[f_*:H^p(Y,\mu_{\ell^n}^{\otimes i})
\longrightarrow H^{p+2d}(X,\mu_{\ell^n}^{\otimes(i+d)})\]
即可（参见第 6 节）。至少当 $Y,X$ 均在某个\Term{scheme}{概形} $S$ 上\Term{smooth}{光滑}时，借助\Term{relative-cohomology}{相对上同调}\Term{purity-theorem}{纯性定理}
\SourcePage{682}
（SGA 4，第 XVI 讲 3.7）就是如此；猜想在 $Y,X$ 均\Term{regular}{正则}时也如此，依据的是绝对\Term{cohomology}{上同调}\Term{purity-theorem}{纯性定理}，它至少已在 $X$ 为\Term{excellent}{优良}且\Term{characteristic}{特征}为零时得到证明（SGA 4，第 XIX 讲 3.2）。当 $X,Y$ 都在复\Term{number-field}{数域}上\Term{smooth}{光滑}时，也可在 $X$ 的复点组成的\Term{locally-compact}{局部紧}\Term{space}{空间} $X(\mathbf C)$ 上使用整\Term{chern-class}{Chern 类}。Atiyah 与 Hirzebruch [1] 在这种情形下证明了 (3.1)，且他们更\Term{generic}{一般}地考虑复\Term{analytic-manifold}{解析流形}的\Term{immersion}{浸入}（使用“朴素”的 $K'$）。由此，利用“\Term{lefschetz-principle}{Lefschetz 原理}”和 SGA 4 的\Term{generic}{一般}结果可知：只要 $Y,X$ 都在\Term{characteristic}{特征}为零的\Term{algebraically-closed-field}{代数闭域}上\Term{smooth}{光滑}，公式 (3.1) 的 $\ell$ 进版本就成立。因此，\Term{cohomology}{上同调}公式 (3.1) 的\Term{effectivity}{有效性}极其可信；不过要注意，目前即使 $X,Y$ 在\Term{characteristic}{特征}为零的域 $k$ 上\Term{smooth}{光滑}，只要 $k$ 不是\Term{algebraically-closed-field}{代数闭域}，它也还没有得到证明。

\textbf{3.4。}也可采用 de Rham 或 Hodge \Term{form}{形式}的\Term{chern-class}{Chern 类}\Term{cohomology}{上同调}理论 [13]。在此，只有 $X$ 并非\Term{characteristic}{特征}为零时才有必要问 (3.1) 是否成立，因为无论如何，我们都知道（第 6 节）这个公式模去\Term{torsion}{挠元}后成立。对于域 $k$ 上\Term{smooth}{光滑}\Term{quasi-projective}{拟射影}\Term{scheme}{概形}，(3.1) 的 de Rham 或 Hodge 版本还将由 (3.1) 的模\Term{algebraic-equivalence}{代数等价}\Term{cycle-class}{圈类}版本\Term{pushout}{推出}。因此，de Rham 或 \Term{hodge-cohomology}{Hodge 上同调}版本又提供一种可行的方式，来检验 (3.1) 的“\Term{cycle-class}{圈类}”版本是否成立。

\textbf{3.5。}(3.1) 的\Term{transcendental-cohomology}{超越上同调}版本的证明 [1] 本质上是\Term{transcendental}{超越}性的，使用了 $Y$ 在 $X$ 中的\Term{tubular-neighbourhood}{管状邻域}。这提示我们：证明 (3.1) 一类的公式，可能要系统考虑\Term{support}{支撑}在 $Y$ 中的 $K'$ 群和\Term{cohomology-group}{上同调群}；这种考虑也许在技术上就足以实现“在 $Y$ 周围\Term{localization}{局部化}”。

\SGAsection{4}{$K'(X)$ 与 Chow 环 $A(X)$ 的关系}
\textbf{4.1。}为简单起见，假定 $X$ 是域 $k$ 上\Term{smooth}{光滑}\Term{quasi-projective}{拟射影}\Term{scheme}{概形}（更\Term{generic}{一般}的情形见第 8 节），于是有
\SourcePage{683}
模\Term{rational-equivalence}{有理等价}的\Term{cycle-class}{圈类}所成的 \Term{chow-ring}{Chow 环} $A(X)$ [4]。这个环本来就能在相当更\Term{generic}{一般}的情形下定义；它的变体包括 $\operatorname{Gr}^\bullet(X)$（$K'(X)$ 的 $\lambda$-\Term{filtration}{滤过}的\Term{associated-graded}{伴随分次环}）和 $\operatorname{Gr}_{\mathrm{top}}^\bullet(X)$（$K'(X)\simeq K_\bullet(X)$ 按\Term{support}{支撑}的\Term{codimension}{余维数}给出的\Term{filtration}{滤过}的\Term{associated-graded}{伴随分次环}）。这三个环中各有相应的\Term{chern-class}{Chern 类}理论（$A(X)$ 的情形见 [9]，$\operatorname{Gr}^\bullet(X)$ 的情形见本讨论班，它又映入 $\operatorname{Gr}_{\mathrm{top}}^\bullet(X)$）。此外还有\Term{chern-class}{Chern 类}的\Term{cohomology}{上同调}理论（例如 $\ell$ 进理论，$\ell$ 与 $k$ 的\Term{characteristic}{特征}互素；SGA 5，第 VII 讲）：
\[c_i(E)\in H^{2i}(X,\mathbf Z_\ell(i)),\]
其中 $E$ 是 $X$ 上的\Term{vector-bundle}{向量丛}，也可以是 $K'(X)$ 的任意元素。这些\Term{cohomology}{上同调}理论借助下列与\Term{multiplication}{乘法}结构相容的\Term{canonical}{典范}\Term{homomorphism}{同态}联系起来：
\[\begin{tikzcd}[column sep=large,row sep=large]
H^{2*}(X,\mathbf Z_\ell(*)) & {}\\
A(X)\arrow[u,"\gamma"]\arrow[r,"\phi"] &
\operatorname{Gr}_{\mathrm{top}}^\bullet(X)\\
{} & \operatorname{Gr}^\bullet(X)\arrow[u,"j"]
\end{tikzcd}\tag{4.1}\]
这里记
\[H^{2*}(X,\mathbf Z_\ell(*))
=\bigoplus_iH^{2i}(X,\mathbf Z_\ell(i)).\]
\Term{homomorphism}{同态} $j$ 来自 $K'(X)$ 的 $\lambda$-\Term{filtration}{滤过}比按\Term{support}{支撑}\Term{codimension}{余维数}定义的\Term{topological-filtration}{拓扑滤过}更细这一事实（第 X 讲）；它模去\Term{torsion}{挠元}后是\Term{isomorphism}{同构}（第 VII 讲 4.11）。\Term{homomorphism}{同态} $\phi$ 的定义例如见 [RRR]；它把 $X$ 的\Term{codimension}{余维数}为 $i$ 的\Term{closed-subscheme}{闭子概形} $Y$ 对应的\Term{cycle-class}{圈类}映为
\[\operatorname{cl}(\mathscr O_Y)\in
\operatorname{Filt}_{\mathrm{top}}^i(K'(X))
\bmod\operatorname{Filt}_{\mathrm{top}}^{i+1}(K'(X)).\]
这是\Term{surjective}{满}\Term{homomorphism}{同态}，下面还会看到（4.2）它模去\Term{torsion}{挠元}后是\Term{isomorphism}{同构}。所以，模去\Term{torsion}{挠元}后，$A(X)$、$\operatorname{Gr}_{\mathrm{top}}^\bullet(X)$、$\operatorname{Gr}^\bullet(X)$ 三个环\Term{canonical}{典范}\Term{isomorphism}{同构}；因此，用其中任一个环（与 $\mathbf Q$ 作\Term{tensor-product}{张量积}）表述的 \Term{riemann-roch-formula}{Riemann--Roch 公式}最终都是\Term{equivalence}{等价}的。本讨论班主要用\Term{graded-ring}{分次环} $\operatorname{Gr}^\bullet(X)$ 代替 \Term{chow-ring}{Chow 环}，它有一个显著优点：可对任何\Term{locally-ringed-topos}{局部赋环拓扑斯}定义，尤其可对任何\Term{scheme}{概形}
\SourcePage{684}
（\Term{regular}{正则}与否均可）定义。但与 $A(X)$、$\operatorname{Gr}_{\mathrm{top}}^\bullet(X)$ 相比，它也有缺点（4.6）：本身不具有对\Term{proper}{固有}\Term{morphism}{态射}的\Term{covariant}{协变}性（即使只考虑域上\Term{smooth}{光滑}\Term{quasi-projective}{拟射影}\Term{scheme}{概形}）；只有与 $\mathbf Q$ 作\Term{tensor-product}{张量积}后才有这种性质。$A(X)$ 相对于另外两个环 $\operatorname{Gr}_{\mathrm{top}}^\bullet(X)$ 和 $\operatorname{Gr}^\bullet(X)$ 的另一优点是：有从 $A(X)$ 到 $H^{2*}(X,\mathbf Z_\ell(*))$ 的\Term{homomorphism}{同态}，使\Term{cohomology}{上同调}\Term{chern-class}{Chern 类}可以解释为取值于 $A(X)$ 的\Term{chern-class}{Chern 类}的像。这个\Term{homomorphism}{同态}\Term{generic}{一般}不经由 $\operatorname{Gr}_{\mathrm{top}}^\bullet(X)$ \Term{resolution}{分解}（4.7）；而\Term{generic}{一般}是否有从 $\operatorname{Gr}^\bullet(X)$ 到 $H^{2*}(X,\mathbf Z_\ell(*))$ 的\Term{homomorphism}{同态}，将\Term{chern-class}{Chern 类}送到\Term{chern-class}{Chern 类}，至少是可疑的（参见第 5 节）。当然，与 $\mathbf Q$ 作\Term{tensor-product}{张量积}后，这些困难便消失，且得到一个\Term{homomorphism}{同态}
\[\operatorname{Gr}_{\mathrm{top}}^\bullet(X)\longrightarrow
H^{2*}(X,\mathbf Q_\ell(*))
=H^{2*}(X,\mathbf Z_\ell(*))\otimes_{\mathbf Z_\ell}\mathbf Q_\ell,\tag{4.2}\]
将\Term{chern-class}{Chern 类}送到\Term{chern-class}{Chern 类}。

\textbf{4.2。}要证明 (4.1) 的\Term{surjective}{满}\Term{homomorphism}{同态} $\phi$ 模去\Term{torsion}{挠元}后是\Term{isomorphism}{同构}，使用由\Term{chern-class}{Chern 类}理论 [9] \Term{pushout}{推出}的 \Term{chern-homomorphism}{Chern 同态}
\[K'(X)\longrightarrow\operatorname{Ch}(A(X)),\]
其中右边是第 V 讲定义的 $A(X)$ 的 \Term{chern-ring}{Chern 环}。容易看出，它与 $K'(X)$ 的\Term{topological-filtration}{拓扑滤过}及 $\operatorname{Ch}(A(X))$ 的自然\Term{filtration}{滤过}相容，因而在\Term{associated-graded}{伴随分次群}上给出\Term{homomorphism}{同态}
\[\psi:\operatorname{Gr}_{\mathrm{top}}^\bullet(X)\longrightarrow A(X).\tag{4.3}\]
我们将证明公式
\[\begin{aligned}
\psi\phi(x)&=(-1)^{i-1}(i-1)!\,x\quad\bmod\text{挠元},&&x\in A^i(X),\\
\phi\psi(x)&=(-1)^{i-1}(i-1)!\,x\quad\bmod\text{挠元},&&x\in\operatorname{Gr}_{\mathrm{top}}^i(X).
\end{aligned}\tag{4.4}\]
由于 $\phi$ \Term{surjective}{满射}，第一式又蕴含第二式。这些公式显然说明 $\phi$ 与 $\psi$ 模去\Term{torsion}{挠元}后都是\Term{isomorphism}{同构}，并且在\Term{degree}{次数} $i$ 上相差因子 $(-1)^{i-1}(i-1)!$ 而互为逆。

要证明 (4.4) 的第一式，显然可设 $x$ 由 $X$ 的一个\Term{irreducible}{不可约}\Term{closed-subscheme}{闭子概形} $Y$ 定义，且 $Y$ 的\Term{codimension}{余维数}为 $i$；此时公式\Term{equivalence}{等价}于
\SourcePage{685}
\[c_i(\mathscr O_Y)=(-1)^{i-1}(i-1)!\,\operatorname{cl}(Y)
\quad\bmod\text{挠元}.\tag{4.5}\]
左边把 $\mathscr O_Y$ 看作 $X$ 上的\Term{coherent}{凝聚}\Term{sheaf}{层}，因而定义 $K'(X)$ 的一个元素；右边 $\operatorname{cl}(Y)$ 表示 $Y$ 在 $A^i(Y)$ 中的类。由于从 $X$ 中去掉\Term{codimension}{余维数} $>i$ 的\Term{closed-subset}{闭子集}不会改变 $A^i(X)$，可归结到 $Y$ \Term{regular}{正则}的情形；若底域 $k$ 是\Term{perfect-field}{完美域}，则 $Y$ 在 $k$ 上\Term{smooth}{光滑}。但在后一情形，公式 (4.5) 正是包含 $Y\to X$ 及 $K'(Y)$ 的元素 $1$ 的 \Term{riemann-roch-formula}{Riemann--Roch 公式}，其中使用了取值于 \Term{chow-ring}{Chow 环} $A(X)$ 的\Term{chern-class}{Chern 类}理论。[2] 已在这里考虑的条件下证明了这个公式，所用方法与本讨论班中利用 $\operatorname{Gr}^\bullet(X)$ 证明类似公式的方法十分接近。因此，至少对\Term{perfect-field}{完美域} $k$，公式 (4.4) 得证，继而 $\phi$ 模去\Term{torsion}{挠元}后为\Term{isomorphism}{同构}这一事实也得证。（这反过来又说明，对于 $k$ 上\Term{smooth}{光滑}\Term{algebra}{代数}\Term{scheme}{概形}之间的\Term{proper}{固有}\Term{morphism}{态射}，“有分母”的 \Term{riemann-roch-formula}{Riemann--Roch 公式}本质上是相同的，无论是像 [2] 那样用 $A(X)$ 书写，还是像本讨论班那样用 $\operatorname{Gr}^\bullet(X)$ 书写！）对 $k$ 不作假设的情形，很容易归结到 $k$ 是\Term{algebraically-closed-field}{代数闭域}的情形，只需利用如下事实：对于 $k$ 的\Term{degree}{次数}为 $n$ 的\Term{finite-extension}{有限扩张} $k'$，自然映射
\[A(X)\xrightarrow{f^*}A(X\otimes_k k')\]
的核被 $n$ \Term{annihilate}{零化}；这可利用满足 $f_*f^*=n\operatorname{id}$ 的迹\Term{homomorphism}{同态} $f_*:A(X\otimes_k k')\to A(X)$ 看出。

\textbf{4.3。}看来颇有可能，(4.4) 按原式成立，而不只是模去\Term{torsion}{挠元}后成立；\Term{equivalence}{等价}地，(4.5) 并非只在模去\Term{torsion}{挠元}后成立。这是\Term{immersion}{浸入} $Y\to X$ 的“无分母 \Term{riemann-roch-formula}{Riemann--Roch 公式}” (3.1) 在 $A(X)$ 中\Term{effectivity}{有效性}问题的一个很特殊的情形，即 $i=d$ 的情形。\footnote{因此，借助 P. Jouanolou 的工作，对于域上\Term{smooth}{光滑}\Term{quasi-projective}{拟射影}的 $X$ 已经成立（参见原讲第 5 页的脚注，即本译文 PDF 原页 681）。}
这个特殊情形的\Term{effectivity}{有效性}可以认为是相当可信的。事实上，第 3 节已指出，当 $k$ 的\Term{characteristic}{特征}为零时，无分母 \Term{riemann-roch-formula}{Riemann--Roch 公式}至少在 $\operatorname{Gr}_{\mathrm{top}}^\bullet(X)$ 中成立；这说明 (4.5) 的两边在 $\operatorname{Gr}_{\mathrm{top}}^\bullet(X)$ 中有相同的像，\Term{equivalence}{等价}地，(4.4) 的第二式成立。一旦第 1 节提出的问题得到满意解决，这个结果就会自动推广到任意\Term{characteristic}{特征}。
\SourcePage{686}
作为 (4.5) 在 $A(X)$ 本身中成立的佐证，指出以下事实：左边\Term{restriction}{限制}到 $X-Y$ 后显然为零，因此假设 $Y$ \Term{irreducible}{不可约}时，左边先验地必为 $\operatorname{cl}(Y)$ 的整数倍：
\[c_i(Y)=n(Y)\operatorname{cl}(Y).\]
当 $\operatorname{cl}(Y)$ 不是 $A^i(X)$ 的\Term{torsion}{挠元}时，这就证明了精确公式 (4.5)。若 $X$ 在域 $k$ 上\Term{projective}{射影}（而不只是\Term{quasi-projective}{拟射影}），这一点立即成立，因为若 $\xi\in A^1(X)$ 是 $X$ 的一个\Term{polarization}{极化}，则 $\deg(Y\cdot\xi^i)>0$；这就证明了此情形所需的公式。因此，只要 $X$ \Term{isomorphism}{同构}于某个 $k$ 上\Term{smooth}{光滑}\Term{projective}{射影}\Term{scheme}{概形}的\Term{open-subscheme}{开子概形}，公式仍成立；而若\Term{resolution-of-singularities}{奇点消解}问题有适当的解，这就是事实。特别地，借助 Hironaka [11]，当 $k$ 的\Term{characteristic}{特征}为零时就是如此。

精确公式 (4.5)，因而 (4.4) 的\Term{effectivity}{有效性}，将蕴含
\[\phi^i:A^i(X)\longrightarrow\operatorname{Gr}_{\mathrm{top}}^i(X)\]
的核被 $(i-1)!$ \Term{annihilate}{零化}，特别是
\[\phi^2:A^2(X)\longrightarrow\operatorname{Gr}_{\mathrm{top}}^2(X)\]
为\Term{isomorphism}{同构}。（$\phi^1$ 为\Term{isomorphism}{同构}已可立即由第 X 讲 5.3.2 得到。）如无差错，即使在 $\dim X=3$ 的特殊情形，目前也尚未证明 $\phi^2$ 是\Term{isomorphism}{同构}。无论如何，笔者在 [9，第 4 节，$3^\circ$] 中宣称 (4.4)、(4.5) 成立，是过于仓促了。

\textbf{4.4。}目前，$\operatorname{Gr}_{\mathrm{top}}^\bullet(X)$（或 $\operatorname{Gr}^\bullet(X)$）相对于 \Term{chow-ring}{Chow 环} $A(X)$ 的主要优点，是在“超额”相交，即\Term{dimension}{维数}过大的相交情形，更适合显式计算某些\Term{cycle}{圈}之间的相交公式。通过 Chow \Term{lemma}{引理}使\Term{cycle}{圈}在其类中充分移动，从而定义相交类的方法，确实不大便于计算。一个鲜明的例子是如下基本公式（第 VII 讲 2.7）：对于 $k$ 上\Term{smooth}{光滑}\Term{quasi-projective}{拟射影}\Term{scheme}{概形}的\Term{closed-immersion}{闭浸入} $f:Y\to X$，有
\[f^*f_*(y)=y\,c_d(\check N).\tag{4.6}\]
这个公式在 $\operatorname{Gr}_{\mathrm{top}}^\bullet(X)$ 中成立，其中 $y\in\operatorname{Gr}_{\mathrm{top}}^\bullet(Y)$，$N$ 是 $Y$ 在 $X$ 中的\Term{conormal-sheaf}{余法层}，
\SourcePage{687}
假定它处处秩为 $d$。在上述出处中，这一公式来自一个很简单的 $\operatorname{Tor}$ \Term{sheaf}{层}计算。由 4.2 可知，上式蕴含 $A(Y)$ 中的同一公式，但只在模去\Term{torsion}{挠元}后成立。令人惊讶的是，目前还没有证明 (4.6) 在 $A(Y)$ 本身中成立，\footnote{参见原讲第 23 页的补注（PDF 原页 699）中标为 $(*)$ 的注。}甚至取 $y=1$ 得到的 $Y$ 的\Term{self-intersection}{自交}这一特殊情形也没有证明：
\[f^*(\operatorname{cl}(Y))=c_d(\check N),\tag{4.7}\]
即使 $k$ 是复\Term{number-field}{数域}也是如此！

试图仿照计算 $\operatorname{Gr}^\bullet(X')_{\mathbf Q}$ 的方法（第 VII 讲 4.8），来计算 $X$ 沿 $Y$ 的\Term{blow-up}{爆破} $X'$ 的 \Term{chow-ring}{Chow 环}时，也遇到类似问题；这一方法不与 $\mathbf Q$ 作\Term{tensor-product}{张量积}就同样可以计算 $\operatorname{Gr}_{\mathrm{top}}^\bullet(X')$。完成这项计算的关键公式（沿用上述出处的记号）是
\[f^*(i_*(y))=j_*\bigl(g^*(y)c_{d-1}(\check F)\bigr).\tag{4.8}\]
借助 $\operatorname{Tor}$ 的计算（第 VII 讲 3.4），目前已在 $\operatorname{Gr}_{\mathrm{top}}^\bullet(X')$ 中建立这一公式（$y\in\operatorname{Gr}_{\mathrm{top}}^\bullet(Y)$），但其在 $A(X')$ 中的\Term{effectivity}{有效性}（$y\in A(Y)$）仍待证明。

还应指出，即使 (4.6)、(4.7)、(4.8) 的 $\ell$ 进类似公式目前也尚未证明；当 $y$ 为\Term{algebra}{代数}\Term{cycle}{圈}的\Term{cohomology-class}{上同调类}时，(4.6)、(4.8) 至少在模去\Term{torsion}{挠元}后成立。\footnote{参见原讲第 23 页的补注（PDF 原页 699）中标为 $(**)$ 的注。}

\textbf{4.5。}本节及以下各节将给出一些例子，其中图 (4.1) 的\Term{homomorphism}{同态} $j$ 或 $\phi$ 不是\Term{isomorphism}{同构}。显然，对给定的 $X$，说
\[j:\operatorname{Gr}^\bullet(X)\longrightarrow\operatorname{Gr}_{\mathrm{top}}^\bullet(X)\]
是\Term{isomorphism}{同构}、\Term{monomorphism}{单态射}或\Term{epimorphism}{满态射}是\Term{equivalence}{等价}的。而且，$j$ 的像是 $\operatorname{Gr}_{\mathrm{top}}^\bullet(X)$ 中由 $X$ 上\Term{locally-free}{局部自由}\Term{sheaf}{层}的\Term{chern-class}{Chern 类}生成的\Term{subring}{子环}。如果 $K'(X)$ 作为 $\mathbf Z$-\Term{algebra}{代数}由 $X$ 上\Term{invertible-sheaf}{可逆层}的类生成，则立即可见，这个\Term{subring}{子环}也是由 $\operatorname{Gr}_{\mathrm{top}}^1(X)$，即 $\operatorname{Gr}_{\mathrm{top}}^\bullet(X)$ 的\Term{degree}{次数} $1$ \Term{set-component}{分量}生成的\Term{subring}{子环}。
\SourcePage{688}
现在取 $X$ 为\Term{flag-variety}{旗簇}
\[X=G/B,\]
其中 $G$ 是\Term{algebraically-closed-field}{代数闭域} $k$ 上\Term{simply-connected}{单连通}\Term{semisimple-group}{半单代数群}，$B$ 是一个 \Term{borel-subgroup}{Borel 子群} [5]。当 $k=\mathbf C$ 时，$X$ 的 Bruhat \Term{cell-decomposition}{胞腔分解} [5，第 13 讲] 蕴含\Term{canonical}{典范}\Term{homomorphism}{同态}
\[K'(X)\longrightarrow K(X(\mathbf C))\qquad(*)\]
是\Term{isomorphism}{同构}，而\Term{canonical}{典范}\Term{homomorphism}{同态}
\[A(X)\longrightarrow H^{2*}(X(\mathbf C),\mathbf Z)\qquad(**)\]
也是\Term{isomorphism}{同构}；这里 $X(\mathbf C)$ 是 $X$ 的复点组成的\Term{compact}{紧}\Term{space}{空间}，配备由 $\mathbf C$ 的\Term{topology}{拓扑}\Term{derive}{导出}的通常\Term{topology}{拓扑}。另一方面，L. Hodgkin [12] 证明，$K(X(\mathbf C))$ 确实由 $X$ 上\Term{invertible-sheaf}{可逆层}的类生成。于是 $K'(X)$ 本身也是如此；利用 $\mathbf Z$ 上的\Term{split-group}{分裂}\Term{semisimple-group}{半单群}（SGA 3）作\Term{specialization}{特殊化}，也不难\Term{pushout}{推出}，同一结果对任意\Term{algebraically-closed-field}{代数闭域} $k$ 仍然成立。此外，$X$ 的\Term{cell-decomposition}{胞腔分解}（或在 $k=\mathbf C$ 时 $(**)$ 的\Term{bijective}{双射}性）容易蕴含 $A(X)$ \Term{torsion-free}{无挠}，所以由 4.2，\Term{homomorphism}{同态} $\phi:A(X)\to\operatorname{Gr}_{\mathrm{top}}^\bullet(X)$ 是\Term{isomorphism}{同构}。由此得到：$j$ 是\Term{isomorphism}{同构}，当且仅当 $A(X)$ 由其\Term{degree}{次数} $1$ 的元素生成；或者，在 $k=\mathbf C$ 时，当且仅当 $X(\mathbf C)$ 的整\Term{cohomology}{上同调}由\Term{degree}{次数} $2$ 的元素生成，即\Term{classifying-space}{分类空间} $B_{G(\mathbf C)}$ \Term{torsion-free}{无挠}。另一方面，已知 [15，第 217 页\Term{theorem}{定理} B、$(1)\Leftrightarrow(2)$ 及 2.5]，这成立当且仅当 $G$ “\Term{torsion-free}{无挠}”，也即 $G$ \Term{isomorphism}{同构}于若干 $\operatorname{SL}(n)$ 与 $\operatorname{Sp}(m)$ 的\Term{direct-product}{直积}。因此，只须取 $G$ 为 $n\geq7$ 的 $\operatorname{SO}(n)$ 或一个例外\Term{simple-group}{单群}，就得到 $j$ 不是\Term{isomorphism}{同构}的例子。

\textbf{4.6。}注意，在这个例子中，“正确”的环无疑是\Term{torsion-free}{无挠}环 $A(X)\simeq\operatorname{Gr}_{\mathrm{top}}^\bullet(X)$：它以 Bruhat \Term{cell-decomposition}{胞腔分解}中的胞腔为基，在 $k=\mathbf C$ 时也\Term{isomorphism}{同构}于 $X(\mathbf C)$ 的整\Term{cohomology}{上同调}；相比之下，$\operatorname{Gr}^\bullet(X)$ 则是一个带有相当病态的\Term{torsion}{挠元}的环。因此，这个例子支持这样的看法：\Term{invariants}{不变量} $\operatorname{Gr}^\bullet(X)$ 只有在
\SourcePage{689}
模去\Term{torsion}{挠元}，即与 $\mathbf Q$ 作\Term{tensor-product}{张量积}之后，才真正有用。同时，这个例子还说明，即使只考虑\Term{smooth}{光滑}\Term{quasi-projective}{拟射影}\Term{scheme}{概形}的\Term{closed-immersion}{闭浸入}，$\operatorname{Gr}^\bullet(X)$ 也只有与 $\mathbf Q$ 作\Term{tensor-product}{张量积}后，才具有对\Term{proper}{固有}\Term{morphism}{态射}的\Term{covariant}{协变}性。事实上，若 $X$ 使得 $j$ 不是\Term{isomorphism}{同构}，则 $X$ 中存在\Term{integral}{整}\Term{closed-subscheme}{闭子概形} $Y$，使 $\operatorname{cl}(Y)\in K'(X)$ 不属于 $\operatorname{Filt}^i(X)$，其中 $i=\operatorname{codim}(Y,X)$。可以取 $Y$ 的\Term{dimension}{维数}最小，从而
\[\operatorname{Filt}_{\mathrm{top}}^{i+1}(X)=\operatorname{Filt}^{i+1}(X).\]
设 $Z$ 是 $Y$ 的\Term{singular-locus}{奇异集}，并记 $X'=X-Z$、$Y'=Y-Z$，则可\Term{pushout}{推出} $\operatorname{cl}(Y')\in K'(X')$ 不属于 $\operatorname{Filt}^i(X')$。这里用到：$K'(X)\to K'(X')$ 是与\Term{augmentation}{增广}相容的\Term{surjective}{满射}，其核为
\[\operatorname{Im}\bigl(K(Z)\to K(X)\bigr)
\subset\operatorname{Filt}_{\mathrm{top}}^{i+1}(X)
=\operatorname{Filt}^{i+1}(X).\]
现在考虑包含
\[f:Y'\longrightarrow X',\]
这是 $k$ 上\Term{smooth}{光滑}\Term{quasi-projective}{拟射影}\Term{scheme}{概形}的\Term{codimension}{余维数}为 $i$ 的\Term{closed-immersion}{闭浸入}，并且
\[f_*:K'(Y')\longrightarrow K'(X')\]
与 $\lambda$-\Term{filtration}{滤过}不相容（相差 $i$ 的位移），因为 $f_*(1)$ 不属于 $\operatorname{Filt}^i(X')$。

\textbf{4.7。}沿用 4.5 的记号。注意，由第 IX 讲可知，上述 Hodgkin \Term{theorem}{定理}\Term{equivalence}{等价}于存在\Term{isomorphism}{同构}
\[K'(G)\simeq\mathbf Z,\qquad\text{因而}\quad
\operatorname{Gr}_{\mathrm{top}}^\bullet(G)\simeq\mathbf Z.\]
\Term{homomorphism}{同态} $K'(G/B)\to K'(G)$ 先验地为\Term{surjective}{满射}，其核是由所有 $\operatorname{cl}(L)-1$ 生成的\Term{ideal}{理想}，其中 $L$ 遍历 $X=G/B$ 上与表示 $B\to\mathbf G_m$ 关联的\Term{invertible-sheaf}{可逆模}（它们实际上给出 $X$ 上的所有\Term{invertible-sheaf}{可逆模}，[5，第 15 讲]）。另一方面，类似论证也表明
\[A(G/B)\longrightarrow A(G)\]
为\Term{surjective}{满射}，其核是由表示 $B\to\mathbf G_m$ 定义的 $A^1(G/B)$ 的元素生成的\Term{ideal}{理想}，因而实际上由整个 $A^1(G/B)$ 生成。\footnote{参见第 IX 讲 3.4、2.2。}这说明，$A(G)$ \Term{reduced}{约化}为 $A^0(X)=\mathbf Z$，当且仅当
\SourcePage{690}
“$G$ \Term{torsion-free}{无挠}”，即 $G$ \Term{isomorphism}{同构}于若干 $\operatorname{SL}(n)$ 与 $\operatorname{Sp}(m)$ 的\Term{direct-product}{直积}。因此，恰在这种情况下，\Term{homomorphism}{同态}
\[\phi:A(G)\longrightarrow\operatorname{Gr}_{\mathrm{top}}^\bullet(G)\]
是\Term{isomorphism}{同构}。所以，若 $G$ 有挠，则其底层\Term{scheme}{概形}给出一个 $\phi$ 不是\Term{bijective}{双射}的例子。

\textbf{4.8。}类似地可以找到例子，使 (4.1) 的\Term{homomorphism}{同态} $\gamma$
\[\gamma:A(G)\longrightarrow H^{2*}(G,\mathbf Z_\ell(*))\tag{4.9}\]
不经由 $\operatorname{Gr}_{\mathrm{top}}^\bullet(G)$ \Term{resolution}{分解}；在此，这\Term{equivalence}{等价}于该\Term{homomorphism}{同态}在\Term{degree}{次数} $>0$ 处不为零。限于 $k=\mathbf C$ 的情形（借助适当的\Term{specialization}{特殊化}\Term{theorem}{定理}，可以这样做），问题便归结为研究\Term{homomorphism}{同态}
\[A(G)\longrightarrow H^{2*}(G(\mathbf C),\mathbf Z),\tag{4.10}\]
它也是由\Term{canonical}{典范}\Term{homomorphism}{同态}
\[A(G/B)\xrightarrow{\sim}H^{2*}((G/B)(\mathbf C),\mathbf Z)
\longrightarrow H^{2*}(G(\mathbf C),\mathbf Z)\]
取商得到的。很容易想到，(4.10) 实际上也许是\Term{injective}{单射}，因为另一方面已知，$G(\mathbf C)$ 的挠\Term{prime-number}{素数}恰好就是\Term{finite-group}{有限群}
\[A^+(G)=\bigoplus_{i>0}A^i(G)\]
的阶的素因子。若果真如此，则只要 $G$ 有挠，(4.10) 就不会经由 $\phi$ \Term{resolution}{分解}；因此，对适当的 $\ell$（准确地说，对整除 $A^+(G)$ 的阶的 $\ell$，即 $G$ 的\Term{torsion-coefficient}{挠系数}），(4.9) 也如此。至少可以断言：在复\Term{number-field}{数域}上的群 $G=G_2$ 的情形，(4.10) 在\Term{degree}{次数} $>0$ 处不为零。准确地说，$A^3(G)$ 中存在一个元素，其在 $H^6(G(\mathbf C),\mathbf Z)$ 中的像的阶为 $2$；由此可知，$\ell=2$ 时，(4.9) 不经由 $\phi$ \Term{resolution}{分解}。这里用到的关于 $G_2$ 的事实，\Term{equivalence}{等价}于说 $H^6((G/B)(\mathbf C),\mathbf Z)$ 中存在一个元素，其在 $H^6(G(\mathbf C),\mathbf Z)$ 中的像不为零；借助复\Term{semisimple-group}{半单代数群}与\Term{compact}{紧} Lie 群之间已知的术语对应，这个事实可在 [3] 中找到，是 J. P. Serre 好意告知我的。

\SourcePage{691}
\SGAsection{5}{$\operatorname{Gr}^\bullet(X)$ 与 $H^{2*}(X,\mathbf Z_\ell(*))$ 的关系}
\textbf{5.1。}4.8 的例子并不排除，对于任意\Term{scheme}{概形} $X$，仍可能找到环\Term{homomorphism}{同态}
\[g^\bullet:\operatorname{Gr}^\bullet(X)\longrightarrow
H^{2*}(X,\mathbf Z_\ell(*)),\]
使对 $X$ 上的任意\Term{locally-free}{局部自由}\Term{module}{模} $E$，$g^i$ 都将第 V 讲的\Term{chern-class}{Chern 类} $c_i(E)$ 变为 SGA 5 第 VII 讲的\Term{cohomology}{上同调}\Term{chern-class}{Chern 类}。当然，这里 $\ell$ 表示与 $X$ 的各\Term{residue-characteristic}{剩余特征}不同的\Term{prime-number}{素数}。由于环 $\operatorname{Gr}^\bullet(X)$ 由各 $c_i(E)$ 生成，可以看出，无论如何，$g^\bullet$ 都由与\Term{chern-class}{Chern 类}相容这一条件唯一确定。

\textbf{5.2。}利用 \Term{chern-homomorphism}{Chern 同态}
\[K'(X)\longrightarrow\operatorname{Ch}(H^{2*}(X,\mathbf Z_\ell(*))),\]
分别就左边的\Term{filtration}{滤过}及右边按\Term{degree}{次数}给出的\Term{filtration}{滤过}取\Term{associated-graded}{伴随分次对象}，便得到一个\Term{graded-group}{分次群}\Term{homomorphism}{同态}（但不是环\Term{homomorphism}{同态}！）
\[\psi:\operatorname{Gr}^\bullet(X)\longrightarrow H^{2*}(X,\mathbf Z_\ell(*)),\]
满足关系
\[\psi(c_i(E))=(-1)^{i-1}(i-1)!\,c_i^\ell(E),\]
其中 $c_i^\ell(E)$ 表示 $\ell$ 进\Term{chern-class}{Chern 类}。于是，对各 $\psi^i$ 作除法，可得到环\Term{homomorphism}{同态}
\[g_{\mathbf Q}:\operatorname{Gr}^\bullet(X)\longrightarrow
H^{2*}(X,\mathbf Q_\ell(*))
=H^{2*}(X,\mathbf Z_\ell(*))\otimes_{\mathbf Z_\ell}\mathbf Q_\ell,\]
将\Term{chern-class}{Chern 类}送到\Term{chern-class}{Chern 类}，且左边的像包含在 $H^{2*}(X,\mathbf Z_\ell(*))$ 的像中。因此，只要将 $H^{2*}(X,\mathbf Z_\ell(*))$ 换成它模\Term{torsion-subgroup}{挠子群}的商，$g^\bullet$ 的存在问题就有肯定回答。要定义假定存在的\Term{homomorphism}{同态} $g$ 的各\Term{set-component}{分量} $g^i$，第一个有疑问的情形是 $g^3$：先验地只能定义 $2g^3=\psi^3$。例如，若能在 $X$ 上找到三个\Term{invertible-sheaf}{可逆层} $L_1,L_2,L_3$，满足
\[(L_1-1)(L_2-1)(L_3-1)=0\quad\text{于 }K'(X),\]
\SourcePage{692}
但又满足
\[c_1(L_1)c_1(L_2)c_1(L_3)\ne0
\quad\text{于 }H^6(X,\mu_2^{\otimes3}),\]
那么 $g^3$ 的存在问题就会有否定回答。

\textbf{5.3。}当 $X$ 在域 $k$ 上\Term{smooth}{光滑}且\Term{quasi-projective}{拟射影}时，也可同样询问，是否存在取值于 \Term{chow-ring}{Chow 环} $A(X)$ 的环\Term{homomorphism}{同态}
\[g:\operatorname{Gr}^\bullet(X)\longrightarrow A(X),\]
将\Term{chern-class}{Chern 类}送到\Term{chern-class}{Chern 类}。对此可以重复上述思考；特别地，这些思考表明，若不能构造 $g^i$ 本身，至少可直接构造 $(-1)^{i-1}(i-1)!g^i$。

\textbf{5.4。}本节结束前指出：若不考虑\Term{scheme}{概形} $X$，而考虑一个\Term{topological-space}{拓扑空间} $X$（例如\Term{compact}{紧}\Term{space}{空间}），也会遇到本质上相同的问题。这时仍可考虑 $K'(X)$ 关于其 $\lambda$-\Term{filtration}{滤过}的\Term{associated-graded}{伴随分次环} $\operatorname{Gr}^\bullet(X)$，并希望找到环\Term{homomorphism}{同态}
\[\operatorname{Gr}^\bullet(X)\longrightarrow H^{2*}(X,\mathbf Z),\]
将\Term{chern-class}{Chern 类}送到\Term{chern-class}{Chern 类}。

\SGAsection{6}{上同调 Riemann--Roch 定理与 Gysin 同态}
\textbf{6.1。}为简单起见，这里只考虑具有\Term{ample}{丰沛}\Term{invertible-sheaf}{可逆模}的\Term{scheme}{概形} $X$；这使我们不必区分朴素的 $K'(X)$ 与另一种 $K'(X)$。先考虑\Term{proper}{固有}\Term{complete-intersection-morphism}{完全交态射} $f:X\to Y$ 的“有分母”\Term{riemann-roch-formula}{Riemann--Roch 公式}。由于对于与 $Y$ 的各\Term{residue-characteristic}{剩余特征}互素的\Term{prime-number}{素数} $\ell$，可以定义\Term{locally-free}{局部自由}\Term{module}{模} $E$ 的 $\ell$ 进\Term{chern-class}{Chern 类}
\[c_i(E)\in H^{2i}(X,\mathbf Z_\ell(i))\]
（SGA 5，第 VII 讲），它对应于一个\Term{homomorphism}{同态}
\[K'(X)\longrightarrow\operatorname{Ch}(H^{2*}(X,\mathbf Z_\ell(*))),\tag{6.1}\]
于是就有必要给出 \Term{riemann-roch-formula}{Riemann--Roch 公式}的一个\Term{cohomology}{上同调}变体。为了使公式有意义，唯一还缺少的组成部分是一个 \Term{gysin-homomorphism}{Gysin 同态}
\SourcePage{693}
\[H^n(X,\mathbf Q_\ell(i))\longrightarrow
H^{n-2d}(Y,\mathbf Q_\ell(i-d)),\tag{6.2}\]
其中 $d$ 为\Term{virtual-relative-dimension}{虚相对维数}。假设这个\Term{homomorphism}{同态}已经定义，再考虑第 VIII 讲所给\Term{form}{形式}、取值于 $\operatorname{Gr}^\bullet(Y)_{\mathbf Q}$ 的 \Term{riemann-roch-formula}{Riemann--Roch 公式}，就会发现：\Term{morphism}{态射} $f$ 的 $\ell$ 进 \Term{riemann-roch-formula}{Riemann--Roch 公式}成立，恰好\Term{equivalence}{等价}于下方方块交换：
\[\begin{tikzcd}[column sep=large,row sep=large]
\operatorname{Gr}^\bullet(X)_{\mathbf Q}\arrow[r]\arrow[d,dashed]
&H^{2*}(X,\mathbf Q_\ell(*))\arrow[d,dashed]\\
\operatorname{Gr}^\bullet(Y)_{\mathbf Q}\arrow[r]
&H^{2*}(Y,\mathbf Q_\ell(*))
\end{tikzcd}\tag{6.3}\]
这里水平\Term{arrow}{箭头}由 (6.1) 取\Term{associated-graded}{伴随分次对象}再作除法而定义，使\Term{chern-class}{Chern 类}送到\Term{chern-class}{Chern 类}；两个竖直\Term{arrow}{箭头}分别是第 VIII 讲的\Term{arrow}{箭头}与\Term{cohomology}{上同调}中的 \Term{gysin-homomorphism}{Gysin 同态}（假设已经给定）。

\textbf{6.2。}注意，若 $Y$ 在域 $k$ 上\Term{smooth}{光滑}，则利用 $k$ 上的\Term{relative-cohomology}{相对上同调}\Term{purity-theorem}{纯性定理}，无论如何都能定义 \Term{gysin-homomorphism}{Gysin 同态}（SGA 5，第 IV 讲）。当 $Y$ \Term{excellent}{优良}、\Term{regular}{正则}且\Term{characteristic}{特征}为零时，借助 SGA 4 第 XIX 讲的绝对\Term{cohomology}{上同调}\Term{purity-theorem}{纯性定理}，同一构造也适用；一旦绝对\Term{cohomology}{上同调}\Term{purity-theorem}{纯性定理}在完全\Term{generic}{一般}的情形下得证，这一构造就会自动适用于所有\Term{excellent}{优良}\Term{regular}{正则}的 $Y$，不必\Term{restriction}{限制}\Term{characteristic}{特征}。最后，当 $X,Y$ 都在同一\Term{scheme}{概形} $S$ 上\Term{smooth}{光滑}、$f$ 是 $S$-\Term{morphism}{态射}时，借助\Term{relative-cohomology}{相对上同调}\Term{purity-theorem}{纯性定理}（SGA 4，第 XVI 讲 3.7），同一构造也适用。应当猜想，在这些可定义 \Term{gysin-homomorphism}{Gysin 同态}的情形中，图 (6.3) 确实交换。目前这个断言唯一得到验证的情形，是 $X,Y$ 都在域 $k$ 上\Term{smooth}{光滑}的情形；它来自与 (6.3) 类似的图
\[\begin{tikzcd}[column sep=large,row sep=large]
A(X)\arrow[r]\arrow[d,dashed]&H^{2*}(X,\mathbf Z_\ell(*))\arrow[d,dashed]\\
A(Y)\arrow[r]&H^{2*}(Y,\mathbf Z_\ell(*))
\end{tikzcd}\tag{6.4}\]
\SourcePage{694}
的\Term{commutativity}{交换性}。这里涉及模\Term{rational-equivalence}{有理等价}的\Term{cycle-class}{圈类}所成的环 $A(X)$、$A(Y)$；此\Term{commutativity}{交换性}已在 SGA 5 第 IV 讲中建立。

\textbf{6.3。}注意，在上述每一种可以定义 \Term{gysin-homomorphism}{Gysin 同态}的情形中，实际得到的构造都适用于以 $\mathbf Z_\ell(*)$ 为系数的\Term{cohomology}{上同调}，而不只是以 $\mathbf Q_\ell(*)$ 为系数的\Term{cohomology}{上同调}。因此，当 $f$ 是\Term{immersion}{浸入}时，就可像第 3 节所说明的那样，提出“无分母”\Term{riemann-roch-formula}{Riemann--Roch 公式}的\Term{effectivity}{有效性}问题。当 $X,Y$ 都在域 $k$ 上\Term{smooth}{光滑}时，(6.4) 的\Term{commutativity}{交换性}表明，只须在 \Term{chow-ring}{Chow 环} $A(Y)$ 中证明无分母公式即可，参见 (3.3)。

\textbf{6.4。}若不对 $Y$ 作\Term{regularity}{正则性}假设，即使 $X,Y$ 都是复\Term{number-field}{数域}上的\Term{finite-type}{有限型}\Term{scheme}{概形}，也看不出如何合理定义\Term{cohomology}{上同调} \Term{gysin-homomorphism}{Gysin 同态} (6.2)，甚至对于\Term{codimension}{余维数}为 $1$ 的\Term{regular-immersion}{正则浸入}也是如此。尽管如此，能否找到一个具有合理性质、特别是使方块 (6.3) 交换的定义？显然，这个问题与以下问题有关：对 $Y$ 中\Term{codimension}{余维数}为 $i$ 的\Term{complete-intersection}{完全交} $X$，定义“它的\Term{cohomology-class}{上同调类}”
\[\gamma(X)\in H^{2i}(Y,\mathbf Z_\ell(i)),\]
使没有\Term{excess-component}{超额分支}的相交变为\Term{cup-product}{杯积}。后一问题至少模去\Term{torsion}{挠元}后，即在 $H^{2*}(Y,\mathbf Q_\ell(*))$ 中，有肯定的解决办法：定义
\[\gamma(X)=\frac{1}{(-1)^{i-1}(i-1)!}\,c_i(\mathscr O_X).\]
但我不知道 $c_i(\mathscr O_X)$ 在 $H^{2i}(Y,\mathbf Z_\ell(i))$ 中是否可被 $(i-1)!$ 整除。特别地，设 $Y$ 是复\Term{number-field}{数域}上的三维\Term{projective}{射影}\Term{scheme}{概形}，$X$ \Term{reduced}{约化}为 $Y$ 的一个（\Term{singular}{奇异}）点，则 $c_3(\mathscr O_X)$ 在 $H^6(Y,\mathbf Z_2(3))$ 中是否可被 $2$ 整除？特别地，这个\Term{chern-class}{Chern 类}在 $H^6(Y,\mathbf Z/2\mathbf Z)$ 中是否为零？沿着同一方向，还可提出以下问题：设 $Y$ 是四维（\Term{singular}{奇异}）\Term{projective}{射影}\Term{scheme}{概形}，$X$ 是一个一维\Term{closed-subscheme}{闭子概形}，并且整体上是三个正\Term{divisor}{除子} $D_1,D_2,D_3$ 的交。考虑它们的模 $2$ \Term{chern-class}{Chern 类}之积：
\[c_1(D_1)c_1(D_2)c_1(D_3)\in H^6(Y,\mathbf Z/2\mathbf Z).\]
\SourcePage{695}
这个 $H^6$ 的元素是否与以 $X$ 为交的各 $D_i$ 的具体选择无关？

\SGAsection{7}{完美复形的 Chern 类}
\textbf{7.1。}若想在不作所考虑的\Term{scheme}{概形}具有\Term{ample}{丰沛}\Term{invertible-sheaf}{可逆模}这一颇为人为的假设下，发展有分母或无分母的\Term{cohomology}{上同调} \Term{riemann-roch-formula}{Riemann--Roch 公式}，就必然要面对以下定义问题：对\Term{scheme}{概形} $X$ 上的\Term{perfect-complex}{完美模复形} $L^\bullet$，定义\Term{chern-class}{Chern 类}
\[c_i(L^\bullet)\in H^{2i}(X,\mathbf Z_\ell(i)),\]
以推广 SGA 5 第 VII 讲中\Term{locally-free}{局部自由}\Term{module}{模}的\Term{chern-class}{Chern 类}。对\Term{generic}{一般}\Term{locally-ringed-topos}{局部赋环拓扑斯}也有同一问题（比较 [10]）；对以\Term{continuous}{连续}复值函数\Term{ringed}{赋环}的\Term{paracompact}{仿紧}\Term{space}{空间}上的\Term{perfect-complex}{完美复形}也是如此，这时采用通常的整\Term{cohomology}{上同调}。在\Term{compact}{紧}\Term{space}{空间}的情形，已知用\Term{perfect-complex}{完美复形}构造的 $K'(X)$，与用\Term{locally-free}{局部自由}\Term{module}{模}构造的“朴素”$K'(X)$ 相同（第 IV 讲），于是直接归结为经典定义。但即使只是\Term{locally-compact}{局部紧}\Term{space}{空间}，这个问题似乎已经并不平凡，因为这种\Term{space}{空间}的整\Term{cohomology}{上同调}不总是其各\Term{compact}{紧}\Term{subset}{子集}\Term{cohomology}{上同调}的\Term{inverse-limit}{逆极限}。

\textbf{7.2。}更准确地说，若 $Y$ 是 $X$ 的\Term{closed-subset}{闭子集}，使 $L^\bullet|_{X-Y}$ 在\Term{derived-category}{导出范畴}意义下为零，也即 $L^\bullet$ 的各\Term{cohomology-sheaf}{上同调层}的\Term{support}{支撑}都在 $Y$ 中，则希望定义\Term{support}{支撑}在 $Y$ 中的\Term{chern-class}{Chern 类}
\[c_i^Y(L^\bullet)\in H_Y^{2i}(X,\mathbf Z_\ell(i)),
\qquad\text{分别为 }H_Y^{2i}(X,\mathbf Z),\]
使之在\Term{canonical}{典范}\Term{homomorphism}{同态}
\[H_Y^{2*}(X,-)\longrightarrow H^{2*}(X,-)\]
下给回 7.1 的“通常”\Term{chern-class}{Chern 类}（原文如此）。
\SourcePage{696}
即使后者归结为\Term{locally-free}{局部自由}\Term{module}{模}的\Term{chern-class}{Chern 类}（$X$ 是具有\Term{ample}{丰沛}\Term{invertible-sheaf}{可逆模}的\Term{scheme}{概形}，或是\Term{compact}{紧}\Term{space}{空间}），这些“带\Term{support}{支撑}”的\Term{chern-class}{Chern 类}也还没有构造出来；恐怕只能与 7.1 所要求的\Term{generic}{一般}\Term{chern-class}{Chern 类}一并构造。\footnote{1968 年 10 月补入。L. Illusie 已在 \Term{hodge-cohomology}{Hodge 上同调}的框架内解决这里提出的问题；参见 [13]。}

带\Term{support}{支撑}的\Term{chern-class}{Chern 类}很可能会在无分母\Term{cohomology}{上同调} \Term{riemann-roch-formula}{Riemann--Roch 公式}的证明中起\Term{action}{作用}（参见 3.5）。另一方面，在该公式中引入\Term{support-family}{支撑族}，应当能借此给出有分母或无分母 \Term{riemann-roch-formula}{Riemann--Roch 公式}的改进版本。

\SGAsection{8}{正则概形的 Chow 环}
问题是：对于具有\Term{ample}{丰沛}\Term{invertible-sheaf}{可逆模}的\Term{noetherian}{诺特}\Term{regular}{正则}\Term{scheme}{概形} $X$，定义\Term{cycle}{圈}的\Term{rational-equivalence}{有理等价}，并研究所得的类环，尤其是它的\Term{functorial}{函子性质}，从而推广 $X$ 在域上\Term{smooth}{光滑}时已知的理论 [4]。关键的技术结果仍将是一个“\Term{moving-lemma}{移动引理}”（moving lemma），保证可以在一个\Term{cycle}{圈}的类中改变该\Term{cycle}{圈}，使它与给定\Term{cycle}{圈}按正确的\Term{dimension}{维数}相交。

所考虑的问题本来就相当自然，尤其是我们已看到（4.5--4.8），只要不限于“模去\Term{torsion}{挠元}”的理论，\Term{chow-ring}{Chow 环} $A(X)$ 的替代品 $\operatorname{Gr}^\bullet(X)$、$\operatorname{Gr}_{\mathrm{top}}^\bullet(X)$ 就很难令人满意。设 $X$ 由域 $k$ 上\Term{smooth}{光滑}\Term{projective}{射影}\Term{algebra}{代数}\Term{scheme}{概形} $Y$ 按如下方式得到：取\Term{projecting-cone}{投影锥} $C$ 的顶点处的\Term{local-scheme}{局部概形}，再去掉原点。假定 $A(X)$ 的确是各个 $U-\{0\}$ 的 \Term{chow-ring}{Chow 环}的\Term{direct-limit}{正向极限}，其中 $U$ 遍历 $C$ 中顶点 $0$ 的\Term{open-neighborhood}{开邻域}，则发现 $A(X)$ \Term{canonical}{典范}\Term{isomorphism}{同构}于\Term{quotient-ring}{商环} $A(Y)/\xi A(Y)$，这里 $\xi\in A^1(Y)$ 由所考虑的 $Y$ 的\Term{projective-immersion}{射影浸入}诱导。这说明，了解 $A(X)$ 这种非经典类型的 \Term{chow-ring}{Chow 环}，\Term{equivalence}{等价}于了解 $A(Y)/\xi A(Y)$；可以把它称为 \Term{chow-ring}{Chow 环} $A(Y)$ 的“\Term{primitive-part}{本原部分}”，其含义受\Term{nonsingular}{非奇异}\Term{projective}{射影}\Term{algebraic-variety}{代数簇}的 Lefschetz--\Term{hodge-cohomology}{Hodge 上同调}理论启发。这个例子说明，具有\Term{isolated-singularity}{孤立奇点}的\Term{local-ring}{局部环}的“\Term{local-scheme}{局部} \Term{chow-ring}{Chow 环}”
\SourcePage{697}
是相当不平凡的几何\Term{invariants}{不变量}，无疑将在这些\Term{local-ring}{局部环}的“算术”研究中发挥\Term{action}{作用}。

当然，一旦对具有\Term{ample}{丰沛}\Term{invertible-sheaf}{可逆模}的\Term{noetherian}{诺特}\Term{regular}{正则}\Term{scheme}{概形}建立了合理的 \Term{chow-ring}{Chow 环}概念，就应将第 4 节中这个环与 $\operatorname{Gr}_{\mathrm{top}}^\bullet(X)$ 之间已建立或猜想的关系推广到这一情形，并用这个环表述第 3 节所提出的“无分母 \Term{riemann-roch-formula}{Riemann--Roch 公式}”问题。

\SGAsection{}{参考文献}
\begin{enumerate}
\item[{[1]}] M. F. Atiyah、F. Hirzebruch，《解析\Term{embedding}{嵌入}的 \Term{riemann-roch-theorem}{Riemann--Roch 定理}》（\textit{The Riemann--Roch theorem for analytic embeddings}），\textit{Topology}，第 1 卷，151--166 页（1962）。
\item[{[2]}] A. Borel、J. P. Serre，《\Term{riemann-roch-theorem}{Riemann--Roch 定理}》（\textit{Le théorème de Riemann--Roch}），\textit{Bull. Soc. Math. France}，第 86 卷，97--136 页（1958）。
\item[{[3]}] R. Bott、H. Samelson，\textit{Amer. Journ. of Math.}，第 80 卷，第 1004 页（1958）。
\item[{[4]}] C. Chevalley，《\Term{rational-equivalence}{有理等价}类 I、II》（\textit{Les classes d'équivalence rationelle, I et II}），巴黎高等师范学校 1958 年讨论班，第 2、3 讲。
\item[{[5]}] C. Chevalley，《\Term{algebra}{代数} Lie 群的分类》（\textit{Classification des groupes de Lie algébriques}），巴黎高等师范学校 1956/58 年讨论班。
\item[{[6]}] A. Dold、D. Puppe，《非\Term{additive-functor}{加性函子}的\Term{homology}{同调}及其应用》（\textit{Homologie nicht-additiver Funktoren. Anwendungen}），\textit{Annales de l'Inst. Fourier}，第 XI 卷，201--312 页（1961）。
\item[{[7]}] H. Grauert，《解析层论中的一个\Term{theorem}{定理}》（\textit{Ein Theorem der analytischen Garbentheorie}），\textit{Pub. Math.} 第 5 号，5--64 页（1960）。
\item[{[8]}] A. Grothendieck，《层的类与 \Term{riemann-roch-theorem}{Riemann--Roch 定理}》（\textit{Classes de faisceaux et théorème de Riemann--Roch}），由高等研究院负责油印的讲义，1957 年（本讲中记作 [RRR]，并作为第 0 章的附录收入本讨论班）。
\item[{[9]}] A. Grothendieck，《Chern 类理论》（\textit{La théorie des classes de Chern}），\textit{Bull. Soc. Math. France}，第 86 卷，137--154 页（1958）。
\end{enumerate}
\SourcePage{698}
\begin{enumerate}
\item[{[10]}] A. Grothendieck，《Chern 类与\Term{discrete}{离散}群的\Term{linear-representation}{线性表示}》（\textit{Classes de Chern et représentations linéaires des groupes discrets}），载《\Term{scheme}{概形}\Term{cohomology}{上同调}十讲》（\textit{Dix exposés sur la cohomologie des schémas}），Masson--North Holland Pub. Cie。
\item[{[11]}] H. Hironaka，《\Term{characteristic-zero}{特征零}域上\Term{algebraic-variety}{代数簇}的\Term{resolution-singularities}{奇点消解}》（\textit{Resolution of singularities of an algebraic variety over a field of characteristic zero}），\textit{Ann. of Math.} 第 79 卷，109--326 页（1964）。
\item[{[12]}] L. Hodgkin，学位论文（即将发表？）。
\item[{[13]}] L. Illusie，正在准备中的工作。
\item[{[14]}] J. P. Serre，《\Term{locally-algebraic}{局部代数}，\Term{multiplicity}{重数}》（\textit{Algèbre Locale, Multiplicités}），\textit{Lecture Notes in Math.} 第 11 号（1965），Springer。
\item[{[15]}] A. Borel，《\Term{connected}{连通}紧 Lie 群的\Term{commutant}{交换子群}与\Term{torsion}{挠性}》（\textit{Sous-groupes commutatifs et torsion des groupes de Lie compacts connexes}），\textit{Tohoku Math. Journal}，第 13 卷，216--240 页（1961）。
\item[{[16]}] P. Deligne，M. Artin、A. Grothendieck、J. L. Verdier 著《\Term{topos}{拓扑斯}理论与\Term{scheme}{概形}的\Term{etale-cohomology}{平展上同调}》（\textit{Théorie des topos et cohomologie étale des schémas}，原文记为 SGA 6）中的第 XVIII 讲。\footnote{1969 年 2 月补入。}
\item[{[17]}] U. Mausin，《对应、母题与幺半变换》（\textit{Correspondances, motifs et transformations monoïdales}，俄文），\textit{Mat. Sborik}，第 77（119）卷，第 4 期，1968 年。\footnote{1969 年 2 月补入。}
\end{enumerate}

\SourcePage{699}
\SGAsection{}{补注}
\textbf{原讲第 1 页的补注。}P. Deligne 于 1968 年写出了 $D^-(X)$ 中的 $\bigwedge^i$ 运算和 $K'(X)$ 中的 $\lambda^i$ 运算的理论。他证明，$K'(X)$ 由此成为一个\Term{pre-lambda-ring}{预 $\lambda$ 环}；若 $\mathscr O_X$ 是某个域上的\Term{algebra}{代数}\Term{sheaf}{层}，则 $K'(X)$ 甚至是一个 $\lambda$-\Term{lambda-ring}{环}。看来，对 Deligne 的方法作进一步改进，就可能去掉对 $\mathscr O_X$ 的这个\Term{restriction}{限制}性假设。

\textbf{原讲第 11 页的补注。}

\textbf{$(*)$。}（1969 年 4 月）D. Mumford 告诉我们，他早在 1959 年 Harvard 的一个讨论班上，就已得到 (9.6) 在 \Term{chow-ring}{Chow 环}中的证明（不忽略\Term{torsion}{挠元}）！读者可在 SGA 5 第 VII 讲 9 中找到这个漂亮结果的证明。显然，这个结果包含下一条 $(**)$ 注中指出的相应\Term{cohomology}{上同调}结果。不过，J. P. Jouanolou 证明的\Term{cohomology}{上同调}结果适用条件更\Term{generic}{一般}：无须\Term{quasi-projective-morphism}{拟射影}\Term{hyperplane}{超平面}假设，且可在任意底上的相对\Term{smooth}{光滑}情形下成立。

\textbf{$(**)$。}（1969 年 1 月）J. P. Jouanolou 刚刚对所有 $y$ 证明了这些公式，而且不忽略\Term{torsion}{挠元}；参见 SGA 5 第 VII 讲第 4.9 节。此前，Deligne 与 Manin [17] 已各自独立地用性质颇不相同的（简单）论证指出：只须对“\Term{algebra}{代数}的”$y$ 建立这些公式，甚至只须在 [16] 中对 $y=1$ 建立公式，就可在\Term{generic}{一般}情形下\Term{pushout}{推出}同一关系。


